\chapter{The electroweak stage}
\label{ch:ew}

Read as interactions (\cref{sec:interactions}), the first rung carries an abelian interaction and the mass of one constituent, and the second
rung the strong and the electroweak interactions of one family. This chapter realises the electroweak stage of the second rung: the family
and its charges, the direction that breaks the Pati--Salam algebra, the electroweak direction that the generator of mass provides, the
fermion and vector masses, the split of the masses within the family, and the obstruction to the sum rule that a restricted split obeys.

Part~I has an electroweak stage at the third rung, on a block of rank $64$ chosen by a projector (\PI{Section~7.6.1}). Its numbers recur
here, at the second rung, on the whole spinor module (\cref{sec:partI}). The normalisations are those of \PI{Section~7.6.1}. A mass kernel is
normalised by unit singular values on its support. The generators of the gauge algebra are orthonormal for the trace form
$\langle X,Y\rangle=\Tr(X^\top Y)/64$, and one coupling $g_A$ multiplies them all; this is an assumption, since the product algebra allows
three couplings, and the conventionally normalised coupling, for which $T_{3L}$ has the eigenvalues $\pm\half$ on the doublets, is
$g=\sqrt8\,g_A$. Vector masses are computed from the Gram matrix
\begin{equation}
\label{eq:gram}
K_{ab}(\Phi)=\Tr\bigl([T_a,\Phi]^\top[T_b,\Phi]\bigr)
\end{equation}
of the orthonormal generators $T_a$ at a vacuum value $\Phi$. Its eigenvalues are the squared masses of the action of \PI{Section~7.4.3}
for $g_A=1$, whose scalar kinetic term uses the unnormalised trace; with the normalised kinetic term of the action of \cref{sec:vectors} the
squared masses are $M^2=Zg_A^2K/64$. At the second rung $\su2_L$ is the half $\half(\id+\omega_\cN)$ of the rotations of $\cN$, $\su2_R$ the
other half, and
\[
Y=T_{3R}+\half(B\!-\!L),\qquad Q=T_{3L}+Y,\qquad B\!-\!L=\tfrac23(H_1+H_2+H_3),
\]
with $H_1,H_2,H_3$ the Cartan generators of \cref{thm:rung2}(iv). These generators are antisymmetric real matrices. The charge of a field
under a generator $X$ is an eigenvalue of the Hermitian operator $\hat X=-iX$ on the eigenspace $\Omega=+i$, and traces of charges are traces
of powers of $\hat X$: for example $\Tr\hat Y^2=\tfrac{10}3$ per Lorentz component on the family, while $\Tr Y^2=-\tfrac{10}3$. The kernels
$\Phi_0$, $K_{1u}$ and $K_{1d}$ are defined in \cref{thm:psstage,thm:fermionmasses}.

\section{The first rung as an interaction}
\label{sec:rung1int}

At the first rung the generator of mass carried into $V^-$ is minus the clock, $G_1\,\omega_S=-\gamma^0$. It selects the clock axis, whose
oriented representative $E_1$ maps to $\Gamma$ by (A1); with this orientation the mass kernel is $\Gamma\omega_S=T$, the Higgs vector at the
clock of \cref{ch:hyp}. In the gauged $\mathfrak g_1=\su2$
its stabiliser is $\uu1_J$, with two massive vectors and one massless one. The global phase $\Omega$ moves $T$ in a direction that the gauge
algebra does not reach, the Dirac phase of \PI{Thm~7.12}, and the kernels of $\mathcal M_1$ fixed by $\uu1_J$ form the complex line through
$T$, spanned by $T$ and $T\Omega=J\omega_S$.
\checks{ew-1}

As a classical gauge theory the first rung therefore carries an abelian interaction with unbroken $\uu1_J$, a massive charged pair of
charge $\pm1$, and a massive fermion doublet of charge $\pm\half$, with the Lagrangian of \cref{sec:coupled}. Its $SU(2)$ completion with the
single fermion doublet has the global anomaly of \cref{prop:onedoublet}; its residual charges and classical spectrum remain meaningful. Its
charge $\half J$ is not the electric charge of the family: at the second rung $\half J=H_1$ is a generator of $\su4$ with the values
$\pm\half$ on every field.

\begin{proposition}[The charge of the first rung at the second]
\label{prop:rung1charge}
At the second rung the charge $\half J$ of the first is $H_1=\half E_2E_3$, a Cartan generator of $\su4$, with the values $\pm\half$ on
every left-handed field, while $Q$ takes the values $0$, $\pm\tfrac13$, $\pm\tfrac23$ and $\pm1$. In the trace form, and in the limit
$v_R\to\infty$ of \cref{thm:vectors}:
\begin{enumerate}[label=(\roman*),leftmargin=2.2em,itemsep=1pt]
\item of the squared norm of $H_1$, two thirds lie along colour, one twelfth along the photon direction $Q$, one twentieth along the $Z$,
which is along $T_{3L}-\tfrac38Q$, and one fifth along the $Z'$;
\item of the squared norm of $Q$, three eighths lie along $T_{3L}$, three eighths along $T_{3R}$, and one twelfth along each of $H_1$, $H_2$
and $H_3$.
\end{enumerate}
\checks{ew-15, numerically}
\end{proposition}

The generator of the first rung thus has the squared overlap $\tfrac1{12}$ with the photon direction of the second, the same as the charges
$H_2$ and $H_3$ of the other two planes of the six temporal directions; the three together make $B\!-\!L$, which carries a quarter of the
squared norm of the photon direction. No rescaling and no change of basis turns the spectrum $\pm\half$ of $\half J$ into that of $Q$. These
are statements about generators in the trace metric. The photon direction of the second rung is a new combination, of which the generator
of the first is one part, and whether the gauge field of the first rung continues into that of the second is not addressed.

\section{The family}

\begin{theorem}[One family of the Standard Model]
\label{thm:family}
\begin{enumerate}[label=(\roman*),leftmargin=2.2em,itemsep=1pt]
\item On the sixteen left-handed Weyl fields of $\R^{64}$, colour, $T_{3L}$ and $Y$ take the values of one family of the Standard Model
with a right-handed neutrino,
\[
(\mathbf3,\mathbf2)_{1/6}\oplus(\bar{\mathbf3},\mathbf1)_{-2/3}\oplus(\bar{\mathbf3},\mathbf1)_{1/3}\oplus(\mathbf1,\mathbf2)_{-1/2}\oplus
(\mathbf1,\mathbf1)_{1}\oplus(\mathbf1,\mathbf1)_{0},
\]
with $B\!-\!L=\tfrac13,-\tfrac13,-\tfrac13,-1,1,1$; a colour triplet with $B\!-\!L=\tfrac13$ is a $\mathbf3$, one with $-\tfrac13$ a
$\bar{\mathbf3}$.
\item On these fields $\Tr\hat Y=\Tr\hat Y^3=\Tr\hat T_{3L}^2\hat Y=\Tr\hat C^2\hat Y=0$ for a colour generator $C$, and
$\Tr\hat Y^2=\tfrac{10}3$ per Lorentz component:
the gravitational, $Y^3$, $SU(2)^2Y$ and $SU(3)^2Y$ anomalies vanish.
\end{enumerate}
\checks{ew-2, ew-3, numerically, with the charges rounded to sixths after a bound on the error}
\end{theorem}

The lepton assignments are the standard electroweak ones of \cite{Glashow61,Weinberg67}; the quark charges follow from the branching of
(i), and the right-handed neutrino is that of the Pati--Salam model \cite{PatiSalam}. Which of the two $\su2$ is $\su2_L$ is
fixed by the orientation that defines the left-handed fields (\cref{thm:ps}(v)); here it is the half on which $\omega_\cN=+1$, which acts on
the $\mathbf4$.

\section{The Pati--Salam direction}

\begin{theorem}[The direction that keeps the Standard Model]
\label{thm:psstage}
\begin{enumerate}[label=(\roman*),leftmargin=2.2em,itemsep=1pt]
\item The kernels of $\mathcal M_2$ invariant under $\su3\oplus\su2_L\oplus\uu1_Y$, for the embedding of \cref{thm:family}, form one complex
line, spanned by $\Phi_0$ and $\Omega\Phi_0$, and it lies in the $(\overline{\mathbf{10}},\mathbf1,\mathbf3)$. No other part of
$\mathcal M_2$ contains such a kernel.
\item The stabiliser of $\Phi_0$ in the gauge algebra is $\su3\oplus\su2_L\oplus\uu1_Y$, of dimension $12$. The eigenvalues of the Gram
matrix \cref{eq:gram}, relative to $|\Phi_0|^2$, are $0$ twelve times, $4$ eight times and $20$ once: six leptoquark vectors with
$Y=\pm\tfrac23$ and the pair $W_R^\pm$ with $Y=\pm1$ at one value, and a neutral $Z'$ at five times that value.
\item $\Phi_0$ has rank four, and its range is the kernel of $Y$ on $\R^{64}$: it is a Majorana mass for the right-handed neutrino and for
nothing else.
\item The orbit of $\Phi_0$ under the $21$ gauge generators and its orbit under these and the two global phases $\Omega$ and $G_2$ have the
same dimension, nine, and $[\Omega,\Phi_0]=[G_2,\Phi_0]=-2\,[T_{3R},\Phi_0]$: the phase rotations of the condensate are gauge rotations, so
no separate Goldstone direction, and no Majoron, arises from these phases at this condensate.
\end{enumerate}
\checks{ew-4, ew-5, ew-6, ew-7, numerically}
\end{theorem}

Statement (i) has a short reason. Among the symmetric products of two of the left-handed fields of \cref{thm:family}, only
$\nu^c\nu^c$ is a singlet of $\su3\oplus\su2_L\oplus\uu1_Y$: every other product carries colour, weak isospin or hypercharge. The invariant
kernels therefore form one complex line, supported on the right-handed neutrino. Once the stage is required to keep the algebra of the
Standard Model, its direction is unique up to a phase; its scale $v_R$ is not fixed. Part~I selects the same direction by potentials
(\PI{Section~7.5}); here it is characterised by what it keeps. Spontaneous parity breaking in left--right
symmetric theories was studied in \cite{SenjanovicMohapatra75}; the right-handed triplet used here has the left--right quantum numbers of
$\Delta_R$ in \cite[eq.~(2.4)]{MohapatraSenjanovic81}, and its $\su4$ representation is that of \cite[eqs.~(3.5)--(3.7)]{DingKingLu2021}.
It is the vacuum $\Phi_0$ of \PI{Thms~7.34 and~7.41}, with the same spectrum $0,4,20$ and the same absence of a Majoron.

\section{The electroweak direction}

By (A3) the mass of the second rung is carried by the generator of mass, carried into $V^-$: $\kappa\,G_2\,\omega_S$.

\begin{proposition}[The generator of mass at the second rung]
\label{prop:ewdir}
\begin{enumerate}[label=(\roman*),leftmargin=2.2em,itemsep=1pt]
\item $G_2\,\omega_S=\pm\gamma^0\omega_\cN$ is an antisymmetric, Lorentz-invariant mass kernel with $(G_2\,\omega_S)^2=-\id$, and it lies in
the $(\mathbf1,\mathbf2,\mathbf2)$.
\item Its stabiliser in the gauge algebra is $\su4\oplus\su2_V$, of dimension $18$, and $T_{3L}+T_{3R}$ lies in $\su2_V$.
\item The kernels of $\mathcal M_2$ fixed by this stabiliser form the complex line through $G_2\,\omega_S$.
\item The image of the clock vector, $\Gamma\omega_S$, has a stabiliser of dimension $16$, and $\Phi_0$ one of dimension $12$.
\item No nonzero kernel of $\mathcal M_2$ has a stabiliser of dimension larger than $18$.
\end{enumerate}
\checks{ew-8, for (i) to (iv), numerically}
\end{proposition}

\begin{proof}[Proof of (v)]
The stabiliser of a sum of components in different isotypic parts is the intersection of their stabilisers, so it suffices to consider one
part at a time. In the $(\mathbf1,\mathbf2,\mathbf2)$ the gauge algebra acts on a complex $2\times2$ matrix $\Phi$ by
$\Phi\mapsto U_L\Phi U_R^\dagger$, and the stabiliser of a nonzero $\Phi$ is $\su4\oplus\su2_V$ when $\Phi^\dagger\Phi$ is a multiple of
$\id$ and $\su4\oplus\uu1$ otherwise. Let $\mathfrak h$ be the stabiliser of a nonzero kernel, with $\dim\mathfrak h\ge18$. Its projection to
$\su4$ has dimension at least $12$. A proper subalgebra $\mathfrak k$ of $\su4$ has dimension less than $12$: otherwise its orthogonal
complement for the Killing form, of dimension at most three, would be $\mathfrak k$-invariant, and the kernel of the action of $\mathfrak k$
on it, of dimension at least nine since $\mathfrak k$ acts there through an orthogonal algebra of dimension at most three, would be a nonzero
ideal of $\su4$ inside $\mathfrak k$, against simplicity. So the projection is $\su4$. Then $\mathfrak h\cap\su4$ is an ideal of $\su4$, and
it is not zero, since otherwise $\mathfrak h$ would embed in $\su2\oplus\su2$, of dimension six. Hence $\mathfrak h\supset\su4$, and a kernel
fixed by $\su4$ lies in the $(\mathbf1,\mathbf2,\mathbf2)$, the only part on which $\su4$ acts trivially (\cref{thm:massfield2}).
\end{proof}

The electroweak direction of (A3) is therefore the direction of largest symmetry among the nonzero values of the mass field, unique up to the
gauge group, the global phase and the scale; it is the vacuum $v_R=0$, $|v_u|=|v_d|$ of \PI{Thm~7.42}. Let a potential of the mass field be
invariant under the gauge algebra and the global phases. Its gradient at a point of this line is fixed by the stabiliser, so by (iii) it
lies in the complex line; the phase invariance removes its angular component within the line, and the gradient is radial. On the
fixed-norm sphere the direction is therefore critical, as in the criterion of Michel \cite{Michel71}; in the full space the radial
derivative can survive. The same holds at the first rung for $T$. Whether the potential has a stationary point at a given radius, or its
minimum there, is not decided. Together with $\Phi_0$ the generator of mass leaves $\su3\oplus\uu1_Q$ (\cref{fig:chain}): in this specified
broken phase the unbroken abelian factor is the electric charge, and its gauge field is the photon direction.

The electroweak direction is a transported element, $G_2\,\omega_S=\sigma_2(T_2)\,\omega_S$. The Pati--Salam direction is not, although it
lies, like every kernel, in $\Cl(6,3)\,\omega_S$ (\cref{thm:massfield2}(i)).

\begin{proposition}[The Majorana line is not transported]
\label{prop:notransport}
Let $\mathcal T=\sigma_2(\Cl(V^+_2))$ be the algebra generated by $(\Gamma;P)$. Every kernel $C\omega_S$ with $C\in\mathcal T$ commutes with
the three diagonal weak rotations $S_i=T_i^L+T_i^R$. No nonzero kernel on the complex line of $\Phi_0$ does. Hence no nonzero kernel on that
line has the form $\sigma_2(a)\,\omega_S$.
\checks{ew-21}
\end{proposition}

\begin{proof}
Each $S_i=\half N_jN_k$ commutes with every generator of $\mathcal T$, since each of $N_j,N_k$ anticommutes with it, and with $\omega_S$, for
the same reason. The line of $\Phi_0$ lies in the $(\overline{\mathbf{10}},\mathbf1,\mathbf3)$, a singlet of $\su2_L$ and a triplet of
$\su2_R$, hence a triplet of the diagonal $\su2_V$, whose only fixed vector is zero; this applies to both real directions of the line.
\end{proof}

The argument excludes a construction of $\Phi_0$ from the transported temporal algebra, even after multiplication by $\omega_S$; it does not
exclude a construction that uses the further spatial generators.

\begin{figure}[t]
\centering
\begin{tikzpicture}[font=\small,>=Stealth,
  alg/.style={draw,rounded corners=3pt,align=center,inner sep=5pt,fill=blue!3}]
\node[alg] (top) at (0,0) {$\su4\oplus\su2_L\oplus\su2_R$\\ \footnotesize dimension $21$};
\node[alg] (l) at (-4.2,-2.2) {$\su3\oplus\su2_L\oplus\uu1_Y$\\ \footnotesize dimension $12$};
\node[alg] (r) at (4.2,-2.2) {$\su4\oplus\su2_V$\\ \footnotesize dimension $18$};
\node[alg] (bot) at (0,-4.4) {$\su3\oplus\uu1_Q$\\ \footnotesize dimension $9$};
\draw[->,thick] (top) -- node[above left,align=right,font=\footnotesize]{$v_R\,\Phi_0\in(\overline{\mathbf{10}},\mathbf1,\mathbf3)$\\ Gram eigenvalues $0^{12},4^8,20^1$} (l);
\draw[->,thick] (top) -- node[above right,align=left,font=\footnotesize]{$\kappa\,G_2\,\omega_S\in(\mathbf1,\mathbf2,\mathbf2)$\\ the generator of mass} (r);
\draw[->,thick] (l) -- node[below left,align=right,font=\footnotesize]{$\kappa\,G_2\,\omega_S$\\ $W^\pm$, $Z$: $m_W^2/m_Z^2\to\frac58$} (bot);
\draw[->,thick] (r) -- node[below right,align=left,font=\footnotesize]{$v_R\,\Phi_0$} (bot);
\end{tikzpicture}
\caption{The breaking of the gauge algebra of the second rung by the two vacuum directions (\cref{thm:psstage,prop:ewdir,thm:vectors}), for
the chosen two-stage background with $v_R\kappa\ne0$. The Pati--Salam direction $\Phi_0$ is the only one that keeps the algebra of the
Standard Model; the generator of mass is the nonzero direction of largest symmetry. Together they leave colour and the electric charge, in
either order. The eigenvalues of the Gram matrix at the first step are relative to $|\Phi_0|^2$, with their multiplicities.}
\label{fig:chain}
\end{figure}

\section{Fermion and vector masses}
\label{sec:vectors}

\begin{theorem}[The masses of the family at the generator of mass]
\label{thm:fermionmasses}
\begin{enumerate}[label=(\roman*),leftmargin=2.2em,itemsep=1pt]
\item The kernels of $\mathcal M_2$ invariant under $\su3\oplus\uu1_Q$ form a space of dimension $12$: two each in the
$(\overline{\mathbf{10}},\mathbf1,\mathbf3)$ and the $(\mathbf{10},\mathbf3,\mathbf1)$, four each in the $(\mathbf{15},\mathbf2,\mathbf2)$ and the
$(\mathbf1,\mathbf2,\mathbf2)$. In each bidoublet the up-type kernels, zero on the sectors $d$ and $e$, and the down-type kernels, zero on $u$
and $\nu$, form spaces of dimension two. The sectors $u$, $d$, $\nu$, $e$ of the charge have real dimensions $24$, $24$, $8$ and $8$.
\item An up-type kernel $K_{1u}$ of the $(\mathbf1,\mathbf2,\mathbf2)$ gives one Dirac entry to the up quarks and the neutrinos, a down-type
kernel $K_{1d}$ one to the down quarks and the electron: $m_u=m_\nu^D$ and $m_d=m_e$.
\item $G_2\,\omega_S$ has up-type and down-type parts of equal norm, and all its $64$ singular values are $1$: with it alone every Dirac pair
of the family has the same entry, and the sixteen Weyl fields form eight Dirac fermions of equal mass.
\item A kernel $K_{15}$ of the $(\mathbf{15},\mathbf2,\mathbf2)$, aligned with $K_{1u}$ on the quarks, is $-3K_{1u}$ on the leptons:
$\Phi=\kappa_1K_{1u}+\kappa_{15}K_{15}$ gives $m_{\rm quark}=|\kappa_1+\kappa_{15}|$ and $m_{\rm lepton}=|\kappa_1-3\kappa_{15}|$.
\item For $\Phi=v_R\Phi_0+\kappa K_{1u}$ the neutral sector has the singular values $\half\bigl(\sqrt{v_R^2+4\kappa^2}\mp v_R\bigr)$, four
times each: the seesaw pair. At $(v_R,v_u,v_d)=(20,3,1)$ the masses are $\half(\sqrt{436}\mp20)$ for the neutral pair, $3$ for the up quarks
and $1$ for the down quarks and the electron, as in \PI{Example~7.36}.
\end{enumerate}
\checks{ew-9, ew-10, ew-11, numerically}
\end{theorem}

The masses in (ii) to (v) are those of the kernel itself, that is, for a unit Yukawa coupling; with a common Yukawa coupling $y$ restored
they are multiplied by $|y|$. These are the statements of \PI{Prop.~7.35} at the second rung: the factor of Georgi and Jarlskog
\cite{GeorgiJarlskog79} and the seesaw \cite{Minkowski77,MohapatraSenjanovic80}. What the second rung adds is (iii): with the mass carried by
the generator of mass alone, as (A3) has it, every Dirac entry of the family is $y\kappa$. With $v_R\Phi_0$ present the two neutral masses
are the singular values of the seesaw block. The masses of a family are not equal, so a split needs values of the mass field in other
neutral directions; (A3) alone does not give them (\cref{sec:split}).

\begin{theorem}[The vector masses at the electroweak stage]
\label{thm:vectors}
Let $\Phi=v_R\Phi_0+v_uK_{1u}+v_dK_{1d}$, with the up-type and down-type kernels
$K_{1u},K_{1d}=\half G_2\,\omega_S\mp(T_{3L}+T_{3R})\,\omega_S$ of \cref{prop:ymvalues}(ii), and put $r=v_R^2$ and $C=64(v_u^2+v_d^2)$.
\begin{enumerate}[label=(\roman*),leftmargin=2.2em,itemsep=1pt]
\item In a trace-orthonormal basis adapted to the charges the Gram matrix \cref{eq:gram} is zero on the eight gluons and $16r$ on the six
leptoquarks; it has two equal charged blocks and one neutral block,
\[
K_{\rm ch}=\begin{pmatrix}C&-128v_uv_d\\-128v_uv_d&C+16r\end{pmatrix},\qquad
K_{\rm n}=\begin{pmatrix}C&-C&0\\-C&C+32r&-16\sqrt6\,r\\0&-16\sqrt6\,r&48r\end{pmatrix},
\]
the latter on the normalised axes of $T_{3L}$, $T_{3R}$ and $B\!-\!L$. The characteristic polynomial of $K_{\rm n}$ is
$\lambda\bigl[\lambda^2-(2C+80r)\lambda+128Cr\bigr]$. For $v_R\ne0$ and $(v_u,v_d)\ne(0,0)$ it has exactly one zero eigenvalue, the photon;
it has two if exactly one of $C$ and $r$ vanishes, and three at $\Phi=0$.
\item At $(v_R,v_u,v_d)=(20,1,\half)$ the eigenvalues of $K$ are
\[
0\ (9),\quad 79.360\ (2),\quad 127.872\ (1),\quad 6400\ (6),\quad 6480.640\ (2),\quad 32032.128\ (1),
\]
the numbers of \PI{Example~7.36}, with $m_W^2/m_Z^2=0.62062$; at $v_R=200$ the ratio differs from $\tfrac58$ by $-4.4\times10^{-5}$. The
stabiliser in the gauge algebra is $\su3\oplus\uu1_Q$, and $Q$ annihilates $\Phi$.
\item As $r\to\infty$ at fixed $v_u$, $v_d$ with $C>0$, the light charged root tends to $C$ and the light neutral root to $8C/5$, so $m_W^2/m_Z^2\to\tfrac58=1-\sin^2\theta_W$,
with
\[
\sin^2\theta_W=\frac{|T_{3L}|^2}{|T_{3L}|^2+|Y|^2}=\frac{1/8}{1/8+5/24}=\frac38 .
\]
\item With the generator of mass, $\Phi=20\,\Phi_0+G_2\,\omega_S$, the ratio is $m_W^2/m_Z^2=0.61349$.
\end{enumerate}
\checks{ew-12, ew-19, numerically: (i) against the eigenvalues at three vacua, (ii) against the printed digits of \PI{Example~7.36}}
\end{theorem}

In the action of \cref{sec:coupled} carried to the second rung,
\[
L=\half\Psi^\top K^\mu D_\mu\Psi-\tfrac y2\,\Psi^\top m\,\Psi-\tfrac14F^a_{\mu\nu}F^{a\mu\nu}+\tfrac Z2\,\langle D_\mu m,D^\mu m\rangle-V(m),
\]
with $D_\mu m=\partial_\mu m+g_A\,[A_\mu,m]$, $m\in\mathcal M_2$ and the trace form $\langle\cdot,\cdot\rangle$, the squared vector masses are $M^2=Zg_A^2K/64=Zg^2K/512$. At the
background $v_R\Phi_0+\kappa\,G_2\,\omega_S$ in the decoupling regime $v_R\gg|\kappa|$, where $v_u=v_d=\kappa$, this gives $M_W^2=\tfrac14Zg^2\kappa^2$ and
$M_Z^2=\tfrac25Zg^2\kappa^2$, the form of the Standard Model with $v=\sqrt Z\,\kappa$, and Dirac entries $y\kappa$; the kernel
$\kappa\,G_2\,\omega_S$ alone leaves $\su4\oplus\su2_V$. The vacuum stabiliser fixes the generator $Q$; its gauge-field combination also
depends on the couplings. For independent canonically normalised couplings $g_L$, $g_R$ and $g_4$ of the three factors,
$1/g_Y^2=1/g_R^2+2/(3g_4^2)$, and equal couplings give $\sin^2\theta_W=\tfrac38$ at the scale where they match. This is the value of
\PI{Thm~7.34}(c) and of unified theories with the canonical normalisation of hypercharge \cite{GeorgiQuinnWeinberg74,PDGGUT}; without that scale, the spectrum, the running and the thresholds it is a
normalisation and not a prediction of the low-energy angle. The deviation of $m_W^2/m_Z^2$ from $\tfrac58$ comes from the mixing of $W$ and
$Z$ with $W_R$ and $Z'$; it is of order $\kappa^2/v_R^2$ and vanishes as $v_R\to\infty$, so that $\rho=1$ at tree level up to such
corrections, as in Part~I.

\section{The uneaten direction}

\begin{theorem}[Broken directions at the second rung]
\label{thm:uneaten}
For $\Phi=v_R\Phi_0+v_uK_{1u}+v_dK_{1d}$, with the $21$ gauge generators, and $23$ with the global phases $\Omega$ and $G_2$:
\begin{center}
\small\setlength{\tabcolsep}{5pt}
\begin{tabular}{@{}lccccc@{}}
\toprule
vacuum & stab.\ in 21 & stab.\ in 23 & massive vectors & 23-orbit & uneaten\\
\midrule
generic, $v_R\ne0$, $v_uv_d\ne0$ & 9 & 10 & 12 & 13 & 1\\
$v_R\ne0$, exactly one of $v_u,v_d$ zero & 9 & 11 & 12 & 12 & 0\\
$v_R\ne0$, $v_u=v_d=0$ & 12 & 14 & 9 & 9 & 0\\
$v_R=0$, $v_uv_d\ne0$, $|v_u|\ne|v_d|$ & 16 & 17 & 5 & 6 & 1\\
$v_R=0$, $|v_u|=|v_d|\ne0$ & 18 & 19 & 3 & 4 & 1\\
$v_R=0$, one of $v_u,v_d$ zero & 16 & 18 & 5 & 5 & 0\\
\bottomrule
\end{tabular}
\end{center}
These are the dimensions of \PI{Thm~7.42}, row by row. The vacuum $v_R\Phi_0+\kappa\,G_2\,\omega_S$ is in the first row, and $G_2\,\omega_S$
alone in the fifth. At the generic vacuum the global charge $4B=-G_2+(B\!-\!L)$, with the values $\tfrac43$, $0$ and $-\tfrac43$, is unbroken,
and the broken charge is $\Omega-G_2$, which is $2$ on the $(\mathbf4,\mathbf2,\mathbf1)$ and $0$ on the $(\bar{\mathbf4},\mathbf1,\mathbf2)$.
With all fields left-handed, the anomaly coefficients of $\Omega-G_2$ are $2$, $4$, $\tfrac43$, $\tfrac{16}3$ and $16$ for $SU(3)^2$,
$SU(2)_L^2$, $Y^2$, $Q^2$ and gravitation, and those of $4B$ are $0$, $2$, $-2$ and $0$ for $SU(3)^2$, $SU(2)_L^2$, $Y^2$ and gravitation.
\checks{ew-13, ew-14, numerically}
\end{theorem}

The orbit count proves that at the generic vacuum one tangent of the orbit under the $23$ generators is not a gauge direction. It becomes a
Goldstone mode only if the potential and the rest of the classical action preserve the two global phases, which gauge covariance does not
require. If they do, the generic vacuum has one broken global direction independent of the gauge orbit; its current has a colour anomaly,
which makes it a candidate axion-like direction \cite{PecceiQuinn77,Weinberg78,Wilczek78}, as $1-\Gamma$ is in \PI{Thm~7.42}. The unbroken
$4B$ is baryon number for the classical vacuum stabiliser, with the electroweak anomaly listed. Gauge invariance alone permits explicit
breaking. On the $(\mathbf1,\mathbf2,\mathbf2)$, a complex bidoublet $F\mapsto U_LFU_R^\dagger$, the gauge-invariant quadratic forms span a
space of dimension three, of which only the norm is invariant under $\Omega-G_2$; the other two are $\Re\det F$ and $\Im\det F$, the forms of
\PI{Prop.~7.43}. The broken phase multiplies $F$ by $e^{iq\alpha}$ and $\det F$ by $e^{2iq\alpha}$, so at $\alpha=0$ a term $\mu^2\Re\det F$ has the slope
$-2q\mu^2\Im\det F$ and the second derivative $-4q^2\mu^2\Re\det F$. At the real diagonal backgrounds of the theorem, where
$\det F=v_uv_d$, it gives the uneaten direction a curvature at tree level when $v_uv_d\ne0$ and none when one of them vanishes; for a general
complex $F$ a mass interpretation also requires a stationary vacuum. Its mass and couplings are not determined by the orbit count, which
also gives no decay constant.
\checks{ew-20}

\section{The split of the masses}
\label{sec:split}

With the generator of mass alone the family has one Dirac entry (\cref{thm:fermionmasses}(iii)). The hypothesis has a second source of values
of the mass field: by (A4) a Yang--Mills generator carried into $V^-$ by $\omega_S$ is a value of the mass field, and at one constituent the
mass field consists of exactly these and the Higgs vectors (\cref{thm:massfield}). At the electroweak stage a vacuum value along a generator
must leave $\su3\oplus\uu1_Q$ unbroken, so the generators that can carry one are the neutral generators, the combinations of $T_{3L}$,
$T_{3R}$ and $B\!-\!L$.

\begin{proposition}[The Yang--Mills values of the neutral generators]
\label{prop:ymvalues}
Write $z_f$ for the value of a neutral kernel on the Dirac pair $f\in\{u,d,\nu,e\}$, relative to $G_2\,\omega_S$ and with $\Omega$ as the
imaginary unit, and $X(f_L)$ for the charge of the left-handed member $f_L$ under a generator $X$.
\begin{enumerate}[label=(\roman*),leftmargin=2.2em,itemsep=1pt]
\item For $X=T_{3L}+T_{3R}$, $B\!-\!L$ and $Q$ the kernel $X\omega_S$ has $z_f=-X(f_L)$, the sign being an orientation convention common to
all neutral generators: the values $(-\half,\half,-\half,\half)$, $(-\tfrac13,-\tfrac13,1,1)$ and $(-\tfrac23,\tfrac13,0,1)$ on
$(u,d,\nu,e)$. The kernel $(T_{3L}-T_{3R})\,\omega_S$ is symmetric and gives no mass. Multiplication by $\Omega$ multiplies the values by $i$.
\item $\half G_2\,\omega_S-(T_{3L}+T_{3R})\,\omega_S$ and $\half G_2\,\omega_S+(T_{3L}+T_{3R})\,\omega_S$ are the up-type and the down-type
kernels $K_{1u}$ and $K_{1d}$ of \cref{thm:fermionmasses}, with unit singular values on $u$ and $\nu$, and on $d$ and $e$: the generator of
mass is their sum, and the Yang--Mills value of $T_{3L}+T_{3R}$ is half their difference.
\item $(T_{3L}+T_{3R})\,\omega_S$ lies in the $(\mathbf1,\mathbf2,\mathbf2)$ and $(B\!-\!L)\,\omega_S$ in the $(\mathbf{15},\mathbf2,\mathbf2)$. The
factor $-3$ of Georgi and Jarlskog is the ratio of the $B\!-\!L$ charges of the leptons and the quarks.
\item With their products by $\Omega$, the kernels $G_2\,\omega_S$, $(T_{3L}+T_{3R})\,\omega_S$ and $(B\!-\!L)\,\omega_S$ span six of the
eight real dimensions of the neutral bidoublet kernels. The other two contain a kernel $D$ with the values $(1,-1,-3,3)$: the part of the
$(\mathbf{15},\mathbf2,\mathbf2)$ in which up and down differ.
\end{enumerate}
\checks{ew-16, numerically}
\end{proposition}

The products by $\Omega$ are not Yang--Mills values, since for $X$ in the gauge algebra $\Omega X$ is symmetric; nor is
$\Omega\,G_2\,\omega_S=-S_2$, the spatial volume of the rung (\cref{thm:massfield2}(iv)), the generator of mass. The map $X\mapsto X\omega_S$
is equivariant under $\su4$ and under the diagonal $\su2_V$, but not under $\su2_L\oplus\su2_R$: $[t,X\omega_S]=[t,X]\,\omega_S+X\,[t,\omega_S]$,
and the generators $\half N_i$ of $\mathfrak s_2$ anticommute with $\omega_S$. For instance $B\!-\!L$ commutes with every weak generator,
while $(B\!-\!L)\,\omega_S$ transforms in a weak bidoublet. The kernels $X\omega_S$ are therefore allowed backgrounds constructed from the
neutral generators, and restricting the vacuum to them is an ansatz.
\checks{ew-22}

\begin{definition}[The restricted split]
\label{def:split}
The restricted electroweak vacuum of the mass field is the generator of mass with its rate together with the Yang--Mills value of a neutral
generator, both carried into $V^-$:
\[
m=\kappa\,G_2\,\omega_S+X\omega_S,\qquad X=x_3\,(T_{3L}+T_{3R})+x_B\,(B\!-\!L),\qquad \kappa,\ x_3,\ x_B\in\R .
\]
A part of $X$ along $T_{3L}-T_{3R}$ drops out.
\end{definition}

\begin{theorem}[The masses of the family in the restricted split]
\label{thm:split}
Let $m$ be the vacuum of \cref{def:split}.
\begin{enumerate}[label=(\roman*),leftmargin=2.2em,itemsep=1pt]
\item The signed Dirac entries are $z_f=\kappa-X(f_L)$:
\[
z_u=\kappa-\tfrac{x_3}2-\tfrac{x_B}3,\qquad z_d=\kappa+\tfrac{x_3}2-\tfrac{x_B}3,\qquad
z_\nu=\kappa-\tfrac{x_3}2+x_B,\qquad z_e=\kappa+\tfrac{x_3}2+x_B ,
\]
and the Dirac masses are $|z_f|$, or $|y\,z_f|$ with a common Yukawa coupling $y$ restored.
\item The signed entries obey the sum rule $z_u-z_d=z_\nu-z_e$. Any three of them are reached, with $\kappa=\tfrac14(2z_u+z_d+z_e)$,
$x_3=z_d-z_u$ and $x_B=\tfrac34(z_e-z_d)$, and the fourth follows: $z_\nu=z_u-z_d+z_e$.
\item With $v_R\Phi_0$ added and a nonzero bidoublet part, the stabiliser in the gauge algebra is $\su3\oplus\uu1_Q$. Among the $23$
generators with the global phases it has dimension $11$ at $x_3=\mp2\kappa$, $x_B=0$, the pure up-type and the pure down-type kernels, where
$m_d=m_e=0$ or $m_u=m_\nu^D=0$ and no direction is left uneaten, and $10$ elsewhere. At the pure up-type vacuum the kernels fixed by this
extended stabiliser span the eight real dimensions of the lines of $\Phi_0$, of the neutral $(\mathbf{10},\mathbf3,\mathbf1)$, and of the
up-type kernels of the $(\mathbf1,\mathbf2,\mathbf2)$ and the $(\mathbf{15},\mathbf2,\mathbf2)$; the gauge stabiliser alone, colour and $Q$,
fixes all twelve neutral directions.
\item For every choice of the signs, $|m_\nu^D-m_u|\le m_d+m_e$, and so also when the values are multiplied by phases through $\Omega$.
\end{enumerate}
\checks{ew-17, ew-18, pt-6, numerically; (iv) follows from (ii) by the triangle inequality}
\end{theorem}

The restricted split thus divides the family in a definite way (\cref{fig:diracpairs}). The generator of mass gives the common value
$\kappa$; the Yang--Mills value of $T_{3L}+T_{3R}$ splits up from down by $\mp\tfrac{x_3}2$; and that of $B\!-\!L$ splits quarks from leptons
in the ratio of their charges: the lepton-to-quark coefficient of this contribution is $-3$, analogous to the Clebsch factor of Georgi and
Jarlskog \cite{GeorgiJarlskog79}, whose eigenvalue relations would need a specified texture and matching scale. Each signed entry is the rate minus the charge of the left-handed
field under one neutral generator. Three real values give four entries, so the Dirac entry of the neutrino is fixed by the other three,
$z_\nu=z_u-z_d+z_e$. The hierarchy of the top quark over the bottom quark and the tau motivates reading the heaviest family, within this
ansatz, as close to $x_3=-2\kappa$, $x_B=0$, the pure up-type kernel, which is the vacuum of larger symmetry in (iii); its down-type masses
then measure the distance from that vacuum, and by (iv) its neutrino Dirac mass lies within $m_d+m_e$ of its up-type mass. This is the
equality of the up-type and the neutrino Dirac masses of Pati--Salam theories with one bidoublet and a common $\su4$-singlet Yukawa coupling
\cite{PatiSalam,MohapatraSenjanovic80}, here with a definite correction; a coupling to the $(\mathbf{15},\mathbf2,\mathbf2)$ can split quark
from lepton entries. At $(\kappa,x_3,x_B)=(1,-1.98,0.004)$ the masses of $u$, $d$, $\nu$ and $e$ are $1.98867$, $0.00867$,
$1.994$ and $0.014$.
\checks{ew-17, ew-18, for the numbers}

\begin{figure}[t]
\centering
\begin{tikzpicture}[font=\small,>=Stealth,x=4.6cm,y=2.1cm]
\draw[->,black!50] (-1.05,-1.5) -- (1.0,-1.5) node[right,font=\footnotesize]{$T_{3L}(f_L)$};
\draw[->,black!50] (-1.05,-1.5) -- (-1.05,0.85) node[above,font=\footnotesize]{$(B\!-\!L)(f_L)$};
\foreach \x/\l in {-0.5/{-\tfrac12},0.5/{\tfrac12}} {\draw[black!50] (\x,-1.53) -- (\x,-1.47); \node[below,font=\footnotesize] at (\x,-1.53) {$\l$};}
\foreach \y/\l in {-1/{-1},0.3333/{\tfrac13}} {\draw[black!50] (-1.08,\y) -- (-1.02,\y); \node[left,font=\footnotesize] at (-1.08,\y) {$\l$};}
\draw[black!20,dashed] (-1.05,0.3333) -- (-0.5,0.3333) (-1.05,-1) -- (-0.5,-1) (-0.5,-1.5) -- (-0.5,-1) (0.5,-1.5) -- (0.5,-1);
\draw[thick,blue!60!black] (-0.5,-1) rectangle (0.5,0.3333);
\foreach \x/\y in {0.5/0.3333,-0.5/0.3333,0.5/-1,-0.5/-1} \fill[blue!60!black] (\x,\y) circle (2.2pt);
\node[above=3pt] at (0.5,0.3333) {$u$: $(1,-\tfrac12,-\tfrac13,1)$};
\node[above=3pt] at (-0.5,0.3333) {$d$: $(1,\tfrac12,-\tfrac13,-1)$};
\node[below=3pt] at (0.5,-1) {$\nu$: $(1,-\tfrac12,1,-3)$};
\node[below=3pt] at (-0.5,-1) {$e$: $(1,\tfrac12,1,3)$};
\node[font=\footnotesize] at (0,-0.33) {$z_u-z_d-z_\nu+z_e=8d$};
\end{tikzpicture}
\caption{The four Dirac pairs of the family at the charges $(T_{3L},B\!-\!L)$ of their left-handed members, with the values of the four real
neutral Dirac directions $\bigl(G_2\,\omega_S,\ (T_{3L}+T_{3R})\,\omega_S,\ (B\!-\!L)\,\omega_S,\ D\bigr)$ on each. The generator of mass is
constant, the Yang--Mills values of $T_{3L}+T_{3R}$ and $B\!-\!L$ are affine in the two charges, and $D=6\,T_{3L}(B\!-\!L)$ is their product.
The alternating sum around the rectangle annihilates the affine functions and picks out $D$: the sum rule of \cref{thm:split} is the absence
of the product term (\cref{thm:splitgen}).}
\label{fig:diracpairs}
\end{figure}

\begin{remark}[The electric charge alone]
The Yang--Mills value of the electric charge alone, $X=\lambda Q$, gives $z_f=\kappa-\lambda Q(f_L)$: the neutrino keeps $\kappa$, and
$z_u=\kappa-\tfrac{2\lambda}3$, $z_d=\kappa+\tfrac\lambda3$, $z_e=\kappa+\lambda$. Eliminating $\kappa$ and $\lambda$ gives
$2z_u=5z_d-3z_e$, so $m_u\le\half(5m_d+3m_e)$ everywhere in this ansatz: the down-type masses cannot both be small against the up-type mass.
The split needs $T_{3L}+T_{3R}$ and $B\!-\!L$ as independent values.
\end{remark}

\section{The sum rule is not protected}
\label{sec:potential}

The vacuum of the restricted split, $v_R\Phi_0$ together with \cref{def:split}, is put in by hand: the hypothesis gives no potential for the
mass field. We now show that no symmetry protects the sum rule of \cref{thm:split}(ii): an invariant quartic drives the mass field out of the
restricted ansatz, by an amount that we compute in closed form. With $S$, $F$, $R$ and $L$ the parts of $m$ in the
$(\mathbf{15},\mathbf2,\mathbf2)$, the $(\mathbf1,\mathbf2,\mathbf2)$, the $(\overline{\mathbf{10}},\mathbf1,\mathbf3)$ and the
$(\mathbf{10},\mathbf3,\mathbf1)$, the quartics
\[
I_1=\Tr\bigl((mm^\top)^2\bigr),\qquad I_2=\Tr(SF^\top RR^\top),\qquad I_3=\Tr(FF^\top RR^\top),\qquad I_4=\Tr(LF^\top RF^\top)
\]
are invariant under the gauge algebra and the global phases, since the projections onto the parts are equivariant and the gauge group acts by
$m\mapsto UmU^\top$ with $U$ orthogonal; $I_1$ is invariant under all of $O(64)$. Gradients are taken for the inner product $\Tr(A^\top B)$,
for which the gradient of $I_1$ is $4\,mm^\top m$. Write $G=G_2\,\omega_S$, $S_3=T_{3L}+T_{3R}$, and $P_f$ for the projector onto the charge
sector $f$. The directions $G$, $S_3\omega_S$, $(B\!-\!L)\,\omega_S$ and $D$ are mutually trace-orthogonal, so the component of a kernel along
$D$ is well defined.

\begin{theorem}[Obstruction to the sum rule]
\label{thm:potential}
\begin{enumerate}[label=(\roman*),leftmargin=2.2em,itemsep=1pt]
\item Let $v_R\ne0$ and let $\Phi=v_R\Phi_0+K$ with $K$ a nonzero kernel of \cref{def:split}, for instance $v_R\Phi_0+\kappa G$ with
$\kappa\ne0$. The gauge stabiliser of $\Phi$ is $\su3\oplus\uu1_Q$, and the kernels it fixes are the twelve neutral directions of
\cref{thm:fermionmasses}(i), $D$ among them. At generic such $\Phi$, including $v_R\Phi_0+\kappa G$, the stabiliser among the $23$ generators
with the global phases has dimension ten and fixes the same twelve directions. The gradient of every gauge-invariant potential at $\Phi$
therefore lies in this space, and no gauge or global symmetry forbids a component along $D$. At $v_R=0$ the stabiliser of $\kappa G$ fixes
only the complex line through $G$ (\cref{prop:ewdir}).
\item Let $m=v_R\Phi_0+K$, where $K=\kappa G+x_3S_3\omega_S+x_B(B\!-\!L)\omega_S+dD$ has the values $z_f$ on the Dirac pairs. Then
\begin{equation}
\label{eq:gradI1}
mm^\top m=v_R^3\,\Phi_0+v_R^2\,z_\nu\,GP_\nu+v_R\,z_\nu^2\,\bigl(2\Phi_0-G\Phi_0G\bigr)+\sum_fz_f^3\,GP_f ,
\end{equation}
where $G\Phi_0G$ is a Majorana kernel of $\nu_L$ in the neutral $(\mathbf{10},\mathbf3,\mathbf1)$, and the neutrino block of the generator of
mass is $GP_\nu=\tfrac18\bigl(G-2S_3\omega_S+3(B\!-\!L)\omega_S-D\bigr)$. The gradient $4\,mm^\top m$ of $I_1$ has the component
\begin{equation}
\label{eq:Dcoef}
c_D=\tfrac12\bigl(z_u^3-z_d^3-z_\nu^3+z_e^3-v_R^2\,z_\nu\bigr)
\end{equation}
along $D$. At a vacuum of \cref{def:split}, where $d=0$ and the sum rule holds,
$c_D=-\tfrac12\bigl[3(z_u-z_d)(z_e-z_d)(z_u+z_e)+v_R^2\,z_\nu\bigr]$; at $v_R\Phi_0+\kappa G$ it is $-\half v_R^2\kappa$.
\item At $\Phi=v_R\Phi_0+\kappa G$ the gradients of $I_2$, $I_3$ and $I_4$ are neutral. Besides components along the condensates already
present, such as the term $2\kappa^2v_R\Phi_0$ of the gradient of $I_3$, which rescales $\Phi_0$, the new directions they induce are: for
$I_2$ the up-type $(\mathbf{15},\mathbf2,\mathbf2)$, with the values $(1,0,-3,0)$, in proportion to $v_R^2\kappa$; for $I_3$ the pure up-type
$(\mathbf1,\mathbf2,\mathbf2)$, with the values $(1,0,1,0)$, in proportion to $v_R^2\kappa$; and for $I_4$ only the neutral
$(\mathbf{10},\mathbf3,\mathbf1)$, in proportion to $v_R\kappa^2$. The tadpoles of $I_1$ and $I_2$ have a part along $D$; those of $I_3$ and
$I_4$ do not.
\item At $v_R=0$ the four gradients at $\kappa G$ lie along $G$.
\end{enumerate}
\checks{pt-1, pt-2, pt-3, pt-4, pt-8, pt-9, pt-11, pt-12; (i) is symmetric criticality \cite{Michel71} applied to the checked stabilisers}
\end{theorem}

\begin{proof}[Proof of (ii)]
$G$ is antisymmetric with all singular values one, so $G^\top G=\id$ and $G^2=-\id$; it is neutral, so it commutes with the projectors $P_f$,
which are spectral projectors of $Q^2$. Each of $G$, $S_3\omega_S$, $(B\!-\!L)\omega_S$ and $D$ is, on each sector, a real multiple of $G$
(\cref{prop:ymvalues,thm:splitgen}), so $K=\sum_fz_fGP_f$, $K^2=-\sum_fz_f^2P_f$ and $K^3=-\sum_fz_f^3GP_f$. By \cref{thm:psstage}(iii),
$\Phi_0$ is antisymmetric with $\Phi_0^\top\Phi_0=P_N$, the projector onto $N=\ker Y$, a subspace of dimension four in the sector $\nu$; so
$\Phi_0^2=-P_N$ and $\Phi_0^3=-\Phi_0$. Within the sector $\nu$, $G$ lies in the $(\mathbf1,\mathbf2,\mathbf2)$ and changes $Y$ by $\pm\half$,
so it maps $N$ into its complement $L$ in the sector and $L$ into $N$. Hence $P_NGP_N=0$, $P_NG+GP_N=GP_\nu$ and $\Phi_0G\Phi_0=0$. Since
$m^\top=-m$, $mm^\top m=-m^3$, and the expansion of $(v_R\Phi_0+K)^3$ with these relations gives \cref{eq:gradI1}; the kernel $G\Phi_0G$ is
supported on $L$.

For \cref{eq:Dcoef} it suffices to work in the real slice. The sectors $u$, $d$ and $e$, of real dimensions $24$, $24$ and $8$, carry
$mm^\top=z_f^2$, and the sector $\nu$, of real dimension $8$, carries the Weyl mass matrix
$\bigl(\begin{smallmatrix}0&z_\nu\\z_\nu&v_R\end{smallmatrix}\bigr)$, whose two singular values each occur four times. Hence
\[
I_1=24z_u^4+24z_d^4+8z_e^4+8z_\nu^4+16v_R^2z_\nu^2+4v_R^4 .
\]
The direction $D$ changes $(z_u,z_d,z_\nu,z_e)$ by $(1,-1,-3,3)$ per unit of $d$, it is trace-orthogonal to $G$, $S_3\omega_S$ and
$(B\!-\!L)\omega_S$, and $\Tr(D^\top D)=24+24+72+72=192$. So $c_D=\tfrac1{192}\,\partial I_1/\partial d$, which is \cref{eq:Dcoef}. The four
real directions have the values $(1,1,1,1)$, $(-\half,\half,-\half,\half)$, $(-\tfrac13,-\tfrac13,1,1)$ and $(1,-1,-3,3)$, so a real neutral
Dirac kernel with the values $w_f$ has the component $\tfrac18(w_u-w_d-w_\nu+w_e)$ along $D$; the values $(0,0,1,0)$ give the stated form of
$GP_\nu$. Under the sum rule $z_\nu=z_u-z_d+z_e$, the identity $x^3+y^3+w^3-(x+y+w)^3=-3(x+y)(y+w)(w+x)$ with $(x,y,w)=(z_u,-z_d,z_e)$
gives the last form.
\end{proof}

What this establishes is a lack of symmetry protection, not the size of the violation. At a vacuum of \cref{def:split} with $v_R\ne0$, $c_D$
vanishes only on a hypersurface of the values. For $v_R$ large against the Dirac entries it is dominated by $-\tfrac12v_R^2z_\nu$, and by the
sum rule $|z_\nu|=|z_u-z_d+z_e|\ge m_u-m_d-m_e$, which for a family like the heaviest is close to $m_u$: there the tadpole along $D$ is of
order $v_R^2m_t$. A potential that contains $I_1$ with a coupling $\lambda$ generically sources a value along $D$: the vacuum moves unless its other
invariants cancel the tadpole, which no symmetry requires. The size of the displacement depends on the curvature of the potential: with a stable, nonsingular
restoring curvature $M^2$ along $D$ and negligible mixing, $d\approx-\lambda c_D/M^2$ to first order; in general the displacement is given by
the inverse Hessian on the physical fluctuations. With a value $d$ the sum rule fails by $8d$, $z_u-z_d-z_\nu+z_e=8d$, and the inequality of
\cref{thm:split}(iv) holds only up to $8|d|$; it can fail outright (\chk{pt-6}). Terms like $I_4$ give $\nu_L$ a Majorana mass of order
$\lambda\,v_R\kappa^2/M_L^2$, the induced value of the type-II seesaw \cite{MohapatraSenjanovic81}, which \cref{def:split} also excludes.

The sum rule is therefore a property of the values that the restricted ansatz supplies, not a consequence of gauge invariance, and a
left--right symmetry that could relate the two triplets is broken by $\Phi_0$ itself. At the stated neutral background the gradient of a gauge-invariant potential lies in the
fixed space of colour and $Q$, so to first order it preserves colour and charge; directions that carry colour or charge could still
destabilise a stationary point, and stability in them is not examined.

\begin{proposition}[What a quartic selects]
\label{prop:select}
\begin{enumerate}[label=(\roman*),leftmargin=2.2em,itemsep=1pt]
\item For every nonzero $m\in\mathcal M_2$,
\[
\frac{\Tr\bigl((mm^\top)^2\bigr)}{\Tr(mm^\top)^2}\le\frac14 ,
\]
with equality exactly when $m$, as a complex symmetric Weyl mass matrix, has rank one.
\item Within the $(\overline{\mathbf{10}},\mathbf1,\mathbf3)$ equality holds exactly on the orbit of $\Phi_0$ up to a nonzero scale: at fixed
nonzero norm, the kernels of real rank four with a stabiliser of dimension $12$.
\item The $(\mathbf{10},\mathbf3,\mathbf1)$ contains the analogous kernels, of real rank four and with the same ratio $\tfrac14$. The connected
gauge group preserves the two triplets separately, so the trace quartic $I_1$ alone does not distinguish a right-handed condensate from a
left-handed one.
\item On the $(\mathbf1,\mathbf2,\mathbf2)$ the ratio lies between $\tfrac1{64}$, at $G_2\,\omega_S$, and $\tfrac1{32}$, at the pure up-type
and down-type kernels; at $v_uK_{1u}+v_dK_{1d}$ it is $(v_u^4+v_d^4)/\bigl(32(v_u^2+v_d^2)^2\bigr)$.
\end{enumerate}
\checks{pt-5, pt-10, numerically}
\end{proposition}

\begin{proof}[Proof of (i) and (ii)]
Write the kernel as a complex symmetric $16\times16$ Weyl mass matrix with singular values $s_i$. Each $s_i$ occurs four times among the
singular values of the real kernel $m$, so $\Tr((mm^\top)^2)=4\sum_is_i^4$ and $\Tr(mm^\top)=4\sum_is_i^2$, and the ratio is
$\tfrac14\sum_is_i^4/(\sum_is_i^2)^2\le\tfrac14$, with equality exactly when one $s_i$ is nonzero. A symmetric tensor of rank one is $w\otimes w$
with $w\in L$. In the right-handed triplet $w$ lies in the $(\bar{\mathbf4},\mathbf1,\mathbf2)$, a $4\times2$ matrix, and $w\otimes w$ lies in
$\mathrm{Sym}^2\bar{\mathbf4}\otimes\mathrm{Sym}^2\mathbf2=(\overline{\mathbf{10}},\mathbf1,\mathbf3)$ exactly when its component in
$\Lambda^2\bar{\mathbf4}\otimes\Lambda^2\mathbf2=(\mathbf6,\mathbf1,\mathbf1)$, whose coefficients are the $2\times2$ minors of $w$, vanishes:
that is, when $w=a\otimes b$. The group $SU(4)\times SU(2)_R$ with the phase acts transitively on such products of fixed norm, and $\Phi_0$ is
one of them.
\end{proof}

A quartic of the $(\overline{\mathbf{10}},\mathbf1,\mathbf3)$ part with the sign that favours the largest ratio thus selects the orbit of
$\Phi_0$ at fixed length within that part: the Majorana mass of a single Weyl field, the right-handed neutrino. On the
$(\mathbf1,\mathbf2,\mathbf2)$ the same sign favours the pure up-type or down-type kernel, and the opposite sign the generator of mass. The parts
have separate invariants, and at fixed component norms the signs of two couplings favour the indicated orbits in the right triplet and in the
bidoublet. They do not fix the relative orientation of the two condensates. For the generator of mass this orientation is not physical:
$G_2\,\omega_S$ is fixed by $\su2_V$, so a rotation of it by $SU(2)_R$ equals one by $SU(2)_L$, which fixes $\Phi_0$, and every relative
orientation is gauge-equivalent to the aligned one. For a rank-one bidoublet it is physical: at fixed $\Phi_0$ a rotation of $K_{1u}$ alone by
$SU(2)_R$ keeps both quartic ratios but, at a generic angle, lowers the joint gauge stabiliser from dimension nine to eight, keeping colour
and losing $Q$. Selecting the neutral up-type alignment relative to the Majorana condensate therefore requires a mixed invariant or an
alignment assumption. With the single trace $I_1$ of all of $m$, the sign that selects a rank-one condensate does not prefer the right
triplet to the left one.
\checks{ew-23} Once $v_R\ne0$, $I_3$ tilts the bidoublet towards the up-type or the down-type kernel by the sign of its coupling.
The ratio $\kappa/v_R$ is not fixed: $\kappa\ll v_R$, which $\rho=1$ needs, is a tuning of the mass parameter of the bidoublet, the hierarchy
problem of left--right and Pati--Salam theories \cite{MohapatraSenjanovic81}. Whether a potential has its global minimum at the full vacuum,
with the values it induces and stable in every direction, is not decided here.

\section{The general split}
\label{sec:splitgeneral}

The potential adds what \cref{def:split} leaves out, a value along $D$. The values of $D$ on the Dirac pairs, $(1,-1,-3,3)$, are six times the
product of the two charges of the left-handed member, $6\,T_{3L}(f_L)\,(B\!-\!L)(f_L)$ (\cref{fig:diracpairs}).

\begin{definition}[The real neutral Dirac sector]
\label{def:splitgen}
The real neutral Dirac sector of the mass field is
\[
m=\kappa\,G_2\,\omega_S+X\omega_S+d\,D,\qquad X=x_3\,(T_{3L}+T_{3R})+x_B\,(B\!-\!L),\qquad \kappa,\ x_3,\ x_B,\ d\in\R .
\]
\end{definition}

This four-parameter family is the real part of the neutral Dirac sector. The full space of kernels that keep colour and the electric charge
has twelve real dimensions: four complex Dirac entries, the Dirac directions with their products by $\Omega$, and two complex Majorana entries,
the line of $\Phi_0$ and the neutral left triplet. If the left triplet takes a vacuum value, the neutral fermion block and the vector masses
must include it.

\begin{theorem}[The general split]
\label{thm:splitgen}
\begin{enumerate}[label=(\roman*),leftmargin=2.2em,itemsep=1pt]
\item Each Dirac pair has the signed entry $z_f=\kappa-X(f_L)+6d\,T_{3L}(f_L)\,(B\!-\!L)(f_L)$: a function of the two charges of the
left-handed field, affine in each.
\item The map $(\kappa,x_3,x_B,d)\mapsto(z_u,z_d,z_\nu,z_e)$ has determinant $\tfrac{32}3$ and the inverse
\[
\begin{gathered}
\kappa=\tfrac18\bigl(3(z_u+z_d)+z_\nu+z_e\bigr),\qquad x_3=\tfrac14\bigl(3(z_d-z_u)+z_e-z_\nu\bigr),\\
x_B=\tfrac38\bigl((z_\nu+z_e)-(z_u+z_d)\bigr),\qquad d=\tfrac18\,(z_u-z_d-z_\nu+z_e).
\end{gathered}
\]
Every signed Dirac spectrum of one family is reached, and the sum rule of \cref{thm:split}(ii) holds exactly when $d=0$.
\end{enumerate}
\checks{pt-7, numerically}
\end{theorem}

We therefore drop the sum rule in the general form. It is a property of the values that the restricted ansatz supplies, not of the theory
with its potential, in which an allowed term generically sources $d$; with $d$ free the four signed Dirac entries of the family are independent, and no relation among
them is claimed. The Dirac masses are $|y\,z_f|$ with a common Yukawa coupling; the signs, or the phases through $\Omega$, matter as soon as a
neutrino Dirac mass is to be inferred from the charged masses. With one family the sum rule would fix only the Dirac entry of the neutrino,
which is not observed, so it contradicts no measured mass by itself; the measured light neutrino mass is a seesaw eigenvalue, through which
the Dirac entry bears on it, of order
$(m_\nu^D)^2/M_R$, and for a family like the heaviest it would put $m_\nu^D$ within $m_b+m_\tau$ of $m_t$, so that a light mass of about
$0.05$~eV would ask for $M_R$ of order $10^{14}$~GeV. With three families it would become the relation $M_\nu^D-M_u=M_e-M_d$ among mass
matrices; since the potential does not protect it and the data do not require it, the general form does without it.

\section{An illustration: the measured masses}
\label{sec:observed}

The general split gives each Dirac pair of a family its own entry. To give each fermion its measured mass it is enough that the three
families, however they arise (\cref{sec:open}), carry an index on which nothing else acts. A value of the mass field is then a sum
$m=\sum_\alpha D_\alpha\otimes F_\alpha$ of the four directions of \cref{def:splitgen} times symmetric family matrices; with the reading of
\cref{thm:families}(i) these are Lorentz-invariant mass kernels, and with diagonal $F_\alpha$ each family has its own four values.

\begin{example}
\label{ex:observed}
Take the values of the Particle Data Group of 2024 \cite{PDG2024}, as listed in \cite[eq.~(2.2)]{AntuschHinzeSaad2025}: the pole masses of the top
quark, $172.4$~GeV from cross sections, and of the charged leptons, and the $\overline{\rm MS}$ masses $m_u$, $m_d$, $m_s$ at $2$~GeV,
$m_c(m_c)$ and $m_b(m_b)$. With all signs positive, a unit Yukawa coupling and no neutrino Dirac mass, \cref{thm:splitgen}(ii) gives:
\begin{center}
\footnotesize\setlength{\tabcolsep}{3.5pt}
\begin{tabular}{@{}cccc|rrrr@{}}
\toprule
family & up-type & down-type & charged lepton & $\kappa$ & $x_3$ & $x_B$ & $d$\\
\midrule
1 & $u$: $2.16$ MeV & $d$: $4.70$ MeV & $e$: $0.51099895$ MeV & $2.63637$ & $2.03275$ & $-2.38088$ & $-0.253625$ MeV\\
2 & $c$: $1.2730$ GeV & $s$: $93.5$ MeV & $\mu$: $105.6583755$ MeV & $0.525645$ & $-0.85821$ & $-0.472816$ & $0.160645$ GeV\\
3 & $t$: $172.4$ GeV & $b$: $4.183$ GeV & $\tau$: $1776.93$ MeV & $66.4407$ & $-125.719$ & $-65.5523$ & $21.2492$ GeV\\
\bottomrule
\end{tabular}
\end{center}
A neutrino Dirac mass $m_\nu^D$ adds $\bigl(\tfrac18,-\tfrac14,\tfrac38,-\tfrac18\bigr)\,m_\nu^D$ to the four values of its family. One value of
the mass field with diagonal family matrices then gives the nine charged fermions their measured masses at once. For the neutrinos we use the two
mass-squared differences measured in oscillations, $\Delta m^2_{21}=7.49\times10^{-5}$~eV$^2$ and $\Delta m^2_{31}=2.513\times10^{-3}$~eV$^2$,
the normal-ordering best fit of NuFIT~6.0 with the Super-Kamiokande atmospheric data (IC24 with SK) \cite[Table~1]{NuFIT6}; oscillation data
constrain mass-squared differences and mixing parameters, and the lightest mass and the heavy Majorana scale are inputs of this example. Choose $v_R\Phi_0\otimes\id$ with a common $v_R=10^{14}$~GeV and a massless lightest neutrino. The seesaw block
$\bigl(\begin{smallmatrix}0&m_\nu^D\\m_\nu^D&v_R\end{smallmatrix}\bigr)$ has the light singular value $m_i$ exactly when
$m_\nu^D=\sqrt{m_i(v_R+m_i)}$, which agrees with $\sqrt{m_iv_R}$ to the digits shown: the Dirac masses $0$, $29.42$ and $70.80$~GeV give the
light masses $0$, $8.654$ and $50.13$~meV, with the measured splittings.
\checks{ms-1, ms-2, ms-3, ms-4, numerically}
\end{example}

Nothing is predicted here: twelve values are fitted to twelve masses, and the Dirac masses of the neutrinos, the common $v_R$, the massless
lightest neutrino and the family index are inputs. The example fits eigenvalues. The inputs are at different scales and in different
schemes, and a relation among them would need the running to a common scale, which is not done; the values of the table would change with
it, and they always exist. With diagonal family matrices the mass matrices of all sectors are aligned and there is no mixing: the quark and
lepton mixings require a misalignment of the corresponding family mass matrices, and they are open. The $W$, $Z$ and Higgs masses take their
values through the free coefficients of the action and of the potential; with $y=1$ the values that give the top quark its mass dominate the
electroweak scale, a relation not examined here.

\section{Comparison with Part~I}
\label{sec:partI}

The electroweak stage of \PI{Sections~7.4--7.7} lies at the third rung, on a block of rank $64$ chosen by a projector, with colour on the
spatial directions of the second and third constituents and the weak directions among the temporal ones; its gauge algebra is the effective
algebra of the block, its electroweak direction has $v_u$ and $v_d$ free, its fermion masses are free (\PI{Prop.~7.35}), and its potential
selects $\Phi_0$ (\PI{Section~7.5}), with the global minimality of the full vacuum left open. Here the stage lies at the second rung, on the
whole module. Colour lies in $\su4=\mathfrak g_2$ on the temporal directions and the weak interaction on the spatial triple of the second
constituent; the gauge algebra is the one that (A2) and (A4) gauge; the electroweak direction is the generator of mass, $|v_u|=|v_d|$, at
which the Dirac entries are equal; the split comes from the Yang--Mills values of the neutral generators and a term in the product of the
charges that an allowed term of the potential generically sources; and at the stated neutral background the gradient of a potential keeps
colour and charge to first order while sourcing the excluded direction. The numbers of
\PI{Prop.~7.35, Example~7.36 and Thms~7.41 and~7.42} are the same in both, because both stages are one Pati--Salam family with the trace form
of a module of dimension $64$, a $(\overline{\mathbf{10}},\mathbf1,\mathbf3)$ and a $(\mathbf1,\mathbf2,\mathbf2)$, evaluated with the same trace
conventions and the same vacuum values; the rates and the mass ratios also depend on the gauge and scalar metrics, the couplings, the
relative alignment and the magnitudes of the vacuum values, which are the same choices in both. What differs is where the structure comes
from: in Part~I a rung is an arena and the family is found inside the third rung by a projector; here a rung is an interaction, and the
family is the second rung.
\checks{ew-5, ew-7, ew-9 to ew-14, for the coincidence of the numbers}
